%% toc: P15 164 164.1,164.11
\section{Relevance axioms and Carnap's relevance measure: derivations}
\label{sec:p15-w02b-p15-024-annex}

Throughout, \(h=P(H)\), \(e=P(E)\) and \(q=P(H\mid E)\) when one
source \(E\) and one target \(H\) are in view. A triple \((h,q,e)\) with
\(0<h,e<1\) and \(0\leq q\leq1\) is \emph{feasible} when the four cells
\(eq\), \(e(1-q)\), \(h-eq\) and \(1-h-e(1-q)\) are nonnegative; these
are the masses of \(EH\), \(E\neg H\), \(\neg EH\) and \(\neg E\neg H\),
so a feasible triple determines the joint table and conversely. For
fixed \(h\in(0,1)\) and \(q\in[0,1]\), every
\(e\in(0,\min\{h/q,(1-h)/(1-q)\}]\cap(0,1)\) is feasible, reading a
quotient with zero denominator as \(+\infty\).

\begin{annexitem}{annex:p15-w02b-p15-024-axiom-status}{thm:p15-w02b-p15-024-axiom-status}{POP-D with symmetry, and POP-D with LIK-D}
\begin{proof}
(1) Let \(0<P(E),P(H_1),P(H_2)<1\). Amount symmetry gives
\(R(E\Yright H_i)_P=R(H_i\Yright E)_P\) for \(i=1,2\). Apply POP-D
with the fixed target \(E\) and the two sources \(H_1,H_2\):
\(R(H_1\Yright E)_P\bowtie R(H_2\Yright E)_P\) exactly when
\(P(E\mid H_1)\bowtie P(E\mid H_2)\). Substituting the two equalities
gives \(R(E\Yright H_1)_P\bowtie R(E\Yright H_2)_P\) exactly when
\(P(E\mid H_1)\bowtie P(E\mid H_2)\), which is LIK-D.

(2) The hypotheses are those of
Theorem~\ref{thm:p15-r43-characterization}, whose second row gives a
strictly increasing \(f_Q\) with \(R(A\Yright B)_P=f_Q(P(B\mid A)/P(B))\)
for all \(A,B\) with \(0<P(A),P(B)<1\). Since
\(P(B\mid A)/P(B)=P(AB)/(P(A)P(B))\) is unchanged by exchanging \(A\)
and \(B\), so is \(R\). If a measure satisfying POP-D has
\(R(A\Yright B)_P\neq R(B\Yright A)_P\) for one such pair, it is not
amount-symmetric, and by what was just shown it cannot satisfy LIK-D.

(3) For fixed \(H\), \(R_D=q-h\) is strictly increasing in \(q\), so D
satisfies POP-D (Theorem~\ref{thm:p15-r43-rankings}). In the model,
\(15\) and \(30\) are the multiples of \(15\), so \(P(A)=2/30=1/15\);
\(15,18,\ldots,42\) are the ten multiples of \(3\), so \(P(B)=1/3\);
\(18,27,36\) are the multiples of \(9\), so \(P(N)=1/10\). Every
multiple of \(15\) or of \(9\) is a multiple of \(3\), so \(AB=A\) and
\(NB=N\), and \(P(B\mid A)=P(B\mid N)=1\). Further
\(P(A\mid B)=2/10\) and \(P(N\mid B)=3/10\), so
\[
 \begin{aligned}
 R_D(B\Yright A)_P&=\tfrac15-\tfrac1{15}=\tfrac2{15},\\
 R_D(B\Yright N)_P&=\tfrac3{10}-\tfrac1{10}=\tfrac15,\\
 R_D(A\Yright B)_P&=1-\tfrac13=\tfrac23 .
 \end{aligned}
\]
LIK-D with source \(B\) and targets \(A,N\) requires equal values,
because the likelihoods \(P(B\mid A)\) and \(P(B\mid N)\) are equal;
D gives \(2/15\neq1/5\). Finally
\(R_Q(B\Yright A)_P=(1/5)/(1/15)=3\) and
\(R_Q(B\Yright N)_P=(3/10)/(1/10)=3\). All events used have
probabilities strictly between zero and one.
\end{proof}
\end{annexitem}

\begin{annexitem}{annex:p15-w02b-p15-024-direction}{thm:p15-w02b-p15-024-direction}{Source rankings in either direction}
\begin{proof}
(1) Let \(R\) be a measure of the stated kind satisfying POP-D. For
the fixed target \(H\) and any two sources \(E_i,E_j\), POP-D gives
\(R(E_i\Yright H)_P\bowtie R(E_j\Yright H)_P\) exactly when
\(P(H\mid E_i)\bowtie P(H\mid E_j)\), for each of
\(\bowtie\in\{<,=,>\}\). So the ranking by \(R\), ties included,
coincides with the ranking by posterior probability, whichever such
\(R\) is chosen. At least one sound ranking exists, since D satisfies
POP-D (item~3 of Theorem~\ref{thm:p15-w02b-p15-024-axiom-status}).
Hence there is exactly one.

(2) All of \(P(H)\), \(P(E_i)\) lie in \((0,1)\), so
Definition~\ref{def:p15-r43-measures} gives
\[
 R_Q(E_i\Yright H)_P=\frac{P(H\mid E_i)}{P(H)}
 =\frac{P(E_iH)}{P(E_i)P(H)}
 =\frac{P(E_i\mid H)}{P(E_i)}=R_Q(H\Yright E_i)_P .
\]
Since \(P(H)>0\) is the same for every \(i\), the map
\(P(H\mid E_i)\mapsto P(H\mid E_i)/P(H)\) is strictly increasing, so the
quotients \(R_Q(E_i\Yright H)_P\), and therefore the reversed quotients,
rank the sources exactly as the posteriors do, which by item~1 is the
sound ranking. If \(R\) satisfies POP-D and LIK-D, item~2 of
Theorem~\ref{thm:p15-w02b-p15-024-axiom-status} makes it
amount-symmetric on events with probabilities in \((0,1)\), so
\(R(H\Yright E_i)_P=R(E_i\Yright H)_P\) for every \(i\), and the
reversed values give the ranking by \(R(E_i\Yright H)_P\), which is
the sound ranking by item~1.

(3) A single model carrying both sources: on the eight atoms fixed by
the truth values of \(H,E_1,E_2\), put \(P(H)=1/2\) and make \(E_1\)
and \(E_2\) independent given \(H\) and given \(\neg H\), with
\[
 P(E_1\mid H)=\tfrac2{25},\quad P(E_1\mid\neg H)=\tfrac1{50},\quad
 P(E_2\mid H)=\tfrac7{10},\quad P(E_2\mid\neg H)=\tfrac3{10}.
\]
Each atom receives \(P(\pm H)\) times a product of two numbers in
\([0,1]\), and the masses sum to one, so this is a probability. Then
\(P(E_1)=\tfrac12(\tfrac2{25}+\tfrac1{50})=\tfrac1{20}\) and
\(P(H\mid E_1)=(\tfrac12\cdot\tfrac2{25})/\tfrac1{20}=\tfrac45\);
\(P(E_2)=\tfrac12(\tfrac7{10}+\tfrac3{10})=\tfrac12\) and
\(P(H\mid E_2)=(\tfrac12\cdot\tfrac7{10})/\tfrac12=\tfrac7{10}\). All
three events have probabilities in \((0,1)\). Since
\(4/5>7/10\), the sound ranking puts \(E_1\) above \(E_2\). For the
reversed D values, Definition~\ref{def:p15-r43-measures} with source
\(H\) and target \(E_i\) gives \(R_D(H\Yright E_i)_P=P(E_i\mid H)-P(E_i)\):
\[
 R_D(H\Yright E_1)_P=\tfrac2{25}-\tfrac1{20}=\tfrac3{100},\qquad
 R_D(H\Yright E_2)_P=\tfrac7{10}-\tfrac12=\tfrac15 ,
\]
which rank \(E_2\) above \(E_1\). The forward values
\(R_D(E_1\Yright H)_P=\tfrac45-\tfrac12=\tfrac3{10}\) and
\(R_D(E_2\Yright H)_P=\tfrac7{10}-\tfrac12=\tfrac15\) give the sound
ranking, as item~1 requires.
\end{proof}
\end{annexitem}

\begin{annexitem}{annex:p15-w02b-p15-024-certain-data}{thm:p15-w02b-p15-024-certain-data}{Likelihoods under certain data}
\begin{proof}
Since \(\neg D H\subseteq\neg D\) and \(P(\neg D)=0\), we have
\(P(\neg D H)=0\) and \(P(DH)=P(H)-P(\neg DH)=P(H)\). For \(P(H)>0\)
this gives \(P(D\mid H)=P(DH)/P(H)=1\). If also \(P(\neg H)>0\), the
same argument with \(\neg H\) gives \(P(D\mid\neg H)=1\). Then
\(0<P(H)<1\) and the denominator likelihood is positive, so
Definition~\ref{def:p15-r43-measures} gives
\(R_L(D\Yright H)_P=1/1=1\).
\end{proof}
\end{annexitem}

\begin{annexitem}{annex:p15-w02b-p15-024-combined-report}{thm:p15-w02b-p15-024-combined-report}{Three conjunction polarities}
\begin{proof}
Write \(x=P(A_3\mid B)\), so \(P(A_3\mid\neg B)=1-x\). The model is
the eight-world law that gives each combination of \(B\) or \(\neg B\),
\(A_1\) or \(\neg A_1\), and \(A_3\) or \(\neg A_3\) the product of
\(P(B)=1/2\) or \(P(\neg B)=1/2\) with the two stated conditional
probabilities (or their complements). Its masses are nonnegative and sum
to one, and \(A_1\), \(A_3\) are independent given \(B\) and given
\(\neg B\) by construction.

By the odds form in Theorem~\ref{thm:p15-r43-measures}, for an event
\(X\) of positive probability,
\(O(B\mid X)=O(B)\,R_L(X\Yright B)_P\), and \(O(B)=1\). So \(X\) is
positive, irrelevant or negative to \(B\) according as
\(R_L(X\Yright B)_P\) is greater than, equal to or less than one.
Here \(R_L(A_1\Yright B)_P=(9/10)/(3/10)=3\), so \(A_1\) is positive;
\(R_L(A_3\Yright B)_P=x/(1-x)\), which is below one exactly when
\(x<1/2\).

Conditional independence gives
\(P(A_1A_3\mid B)=\tfrac9{10}x\) and
\(P(A_1A_3\mid\neg B)=\tfrac3{10}(1-x)\), both positive for
\(0<x<1\), so \(R_L(A_1A_3\Yright B)_P=3x/(1-x)\) and
\(P(B\mid A_1A_3)=3x/(3x+1-x)=3x/(1+2x)\). For \(x=2/5,1/4,1/5\):
\(x/(1-x)=2/3,1/3,1/4\); \(3x/(1-x)=2,1,3/4\); and
\(3x/(1+2x)=(6/5)/(9/5)=2/3\), \((3/4)/(3/2)=1/2\),
\((3/5)/(7/5)=3/7\). Comparing these posteriors with \(P(B)=1/2\) gives
positive, irrelevant and negative relevance respectively, while in all
three models \(A_1\) is positive and \(A_3\) negative to \(B\).
\end{proof}
\end{annexitem}

\begin{annexitem}{annex:p15-w02b-p15-024-carnap-bridge}{thm:p15-w02b-p15-024-carnap-bridge}{From Carnap's measure to Merin's covariance}
\begin{proof}
(1) By definition \(P_e(x)=m(xe)/m(e)\), so
\[
 \begin{aligned}
 \Delta_{P_e}(i,h)&=\frac{m(eih)}{m(e)}-\frac{m(ei)\,m(eh)}{m(e)^2}\\
 &=\frac{m(eih)\,m(e)-m(eh)\,m(ei)}{m(e)^2}
 =\frac{r(i,h,e)}{m(e)^2},
 \end{aligned}
\]
the last step by Carnap's definition D67-1 as recorded in
Definition~\ref{def:p15-carnap-relevance}. For tautological \(e\),
\(m(e)=1\), \(P_e=m\), and \(r(i,h,e)=m(ih)-m(i)m(h)=R_C(i\Yright h)_m\).

(2) Suppose first \(m(ei)>0\). By
Theorem~\ref{thm:p15-carnap-relevance-core},
\(c(h,ei)-c(h,e)=r(i,h,e)/(m(e)m(ei))\) with a positive denominator, so
D65-1a, b and d(1) classify \(i\) by the sign of \(r(i,h,e)\). By (1),
that is the sign of \(\Delta_{P_e}(i,h)\), which by
Definition~\ref{def:p15-r43-polarity} gives \(\mathrm{pos}_{P_e}\),
\(\perp_{P_e}\) or \(\mathrm{neg}_{P_e}\). Suppose next \(m(ei)=0\).
In a finite system a regular measure is positive on every sentence that
is not L-false, so \(e\land i\) is L-false and D65-1d(2) makes \(i\)
irrelevant. On Merin's side \(P_e(i)=0\), hence \(P_e(ih)=0\) and
\(\Delta_{P_e}(i,h)=0\), so \(i\perp_{P_e}h\). The cases agree.

For the additivity laws, Theorem~\ref{thm:p15-carnap-relevance-core}
gives \(r(i\lor j,h,e)=r(i,h,e)+r(j,h,e)\) when \(m(eij)=0\), and
\(r(ij,h,e)=r(i,h,e)+r(j,h,e)\) when \(m(e\neg i\neg j)=0\). Dividing
by \(m(e)^2\) and using (1) turns these into the same identities for
\(\Delta_{P_e}\). The conditions become \(P_e(ij)=0\) and
\(P_e(\neg i\neg j)=0\), that is \(P_e(i\lor j)=1\). These are the
disjoint-source and exhaustive-source laws for \(R_C\) in
Theorem~\ref{thm:p15-r43-composition}, applied to \(P_e\).

(3) First the characterization for an arbitrary probability on a finite
set. Suppose \(A\,\mathrm{expos}_P\,B\): \(P(A)>0\), \(P(B)<1\) and
\(P(B\mid A)=1\). Then \(P(AB)=P(A)>0\), so
\(P(A\neg B)=P(A)-P(AB)=0\), and
\(P(\neg A\neg B)=P(\neg B)-P(A\neg B)=1-P(B)>0\). Conversely, if
\(P(A\neg B)=0\), \(P(AB)>0\) and \(P(\neg A\neg B)>0\), then
\(P(A)=P(AB)+P(A\neg B)=P(AB)>0\), \(P(B\mid A)=P(AB)/P(A)=1\), and
\(P(\neg B)=P(A\neg B)+P(\neg A\neg B)>0\), so \(P(B)<1\).

For tautological \(e\), the four cells of
Theorem~\ref{thm:p15-carnap-relevance-core} are
\(a=m(ih)\), \(b=m(i\neg h)\), \(c=m(\neg ih)\), \(d=m(\neg i\neg h)\).
Theorem~\ref{thm:p15-carnap-relevance-analysis} states that \(i\) is
extremely positive to \(h\) in the sense of D74-1 exactly when \(b=0\),
\(a>0\) and \(d>0\). With \(P=m\), these are the three conditions just
shown equivalent to \(i\,\mathrm{expos}_m\,h\).

In the example, \(P(A\neg B)=P(\{2\})=0\), \(P(AB)=P(\{1\})=1/2\) and
\(P(\neg A\neg B)=P(\{3\})=1/2\), so \(A\,\mathrm{expos}_P\,B\). But
\(2\in A\setminus B\), so \(A\) does not entail \(B\); the nonempty event
\(\{2\}\) has mass zero, so \(P\) is not regular.

(4) Let \(P(H_1H_2)=0\) and \(P(E)>0\). Since \(EH_1H_2\subseteq H_1H_2\),
\(P(E(H_1\lor H_2))=P(EH_1)+P(EH_2)\), and
\(P(H_1\lor H_2)=P(H_1)+P(H_2)\). Hence
\(R_C(E\Yright H_1\lor H_2)_P=R_C(E\Yright H_1)_P+R_C(E\Yright H_2)_P\),
and \(R_C(E\Yright H_1\lor H_2)_P\bowtie R_C(E\Yright H_1)_P\) exactly
when \(\Delta_P(E,H_2)\bowtie0\). Dividing by \(P(E)>0\), this holds
exactly when \(P(H_2\mid E)\bowtie P(H_2)\), which is MDHD. The failure of
POP-D is item~3 of Theorem~\ref{thm:p15-w02b-p15-024-z-characterization}.
\end{proof}
\end{annexitem}

\begin{annexitem}{annex:p15-w02b-p15-024-raven-imbalance}{thm:p15-w02b-p15-024-raven-imbalance}{The raven family}
\begin{proof}
Given \(H\), the four kinds of object (raven and black, raven and
nonblack, nonraven and nonblack, nonraven and black) have masses
\(\rho\), \(0\), \(\nu\), \(1-\rho-\nu\); given \(\neg H\) they have
\(\rho-\varepsilon\), \(\varepsilon\), \(\nu-\varepsilon\),
\(1-\rho-\nu+\varepsilon\). The hypotheses on \(\rho,\nu,\varepsilon\)
make all eight numbers nonnegative, each list sums to one, and the
raven and nonblack marginals are \(\rho\) and \(\nu\) in both strata.
Weighting by \(P(H)\) and \(1-P(H)\) gives a finite law in which no
nonblack raven exists under \(H\), as the universal hypothesis requires.

Thus \(P(A\mid H)=\rho\), \(P(A\mid\neg H)=\rho-\varepsilon>0\),
\(P(B\mid H)=\nu\) and \(P(B\mid\neg H)=\nu-\varepsilon>0\), which gives
the two displayed likelihood ratios, both greater than one. By the odds
form \(O(H\mid X)=O(H)\,R_L(X\Yright H)_P\) of
Theorem~\ref{thm:p15-r43-measures}, both \(A\) and \(B\) raise the
probability of \(H\). With positive denominators,
\[
 \frac{\rho}{\rho-\varepsilon}\bowtie\frac{\nu}{\nu-\varepsilon}
 \iff \rho(\nu-\varepsilon)\bowtie\nu(\rho-\varepsilon)
 \iff -\rho\varepsilon\bowtie-\nu\varepsilon
 \iff \nu\bowtie\rho ,
\]
using \(\varepsilon>0\). Because the prior odds are common and
\(x\mapsto x/(1+x)\) is strictly increasing, \(P(H\mid A)\bowtie
P(H\mid B)\) exactly when the likelihood ratios compare the same way.
By Theorem~\ref{thm:p15-r43-rankings}, every measure satisfying POP-D
ranks the two sources for the fixed target \(H\) by these posteriors.
So \(A\) is more, equally or less relevant according as \(\rho<\nu\),
\(\rho=\nu\) or \(\rho>\nu\). The values \(\rho=1/100\), \(\nu=4/5\),
\(\varepsilon=1/200\) give \(2\) and \(160/159\), the figures of
Section~\ref{sec:p15-r43-rankings}.
\end{proof}
\end{annexitem}

\begin{annexitem}{annex:p15-w02b-p15-024-z-characterization}{thm:p15-w02b-p15-024-z-characterization}{Relative distance: the full realization argument}
\begin{proof}
By formality, \(R(E\Yright H)_P=F(h,q,e)\) for one function \(F\) on
feasible triples. Two constructions are used.

\emph{One target, two sources.} Let \(h\in(0,1)\) and let
\((h,q_1,e_1)\), \((h,q_2,e_2)\) be feasible. Inside \(H\) give \(E_i\)
the probability \(e_iq_i/h\), and inside \(\neg H\) the probability
\(e_i(1-q_i)/(1-h)\); feasibility puts both in \([0,1]\). Make \(E_1\),
\(E_2\) independent within each stratum by multiplying these Bernoulli
probabilities. Then \(P(E_i)=e_iq_i+e_i(1-q_i)=e_i\) and
\(P(HE_i)=e_iq_i\), so \(P(H\mid E_i)=q_i\).

\emph{One source, two targets.} Let \(e\in(0,1)\) and let
\((h_1,q_1,e)\), \((h_2,q_2,e)\) be feasible. Inside \(E\) make
\(H_1,H_2\) independent with probabilities \(q_1,q_2\); outside \(E\)
make them independent with probabilities \((h_i-eq_i)/(1-e)\), which lie
in \([0,1]\) because the cells \(\neg EH_i\) and \(\neg E\neg H_i\) are
nonnegative. Then \(P(H_i)=h_i\) and \(P(H_i\mid E)=q_i\).

\emph{Step 1: the evidence mass drops out.} If \((h,q,e_1)\) and
\((h,q,e_2)\) are feasible, the first construction realizes both in one
model with equal posteriors, and the equality case of POP-D gives
\(F(h,q,e_1)=F(h,q,e_2)\). Since every \((h,q)\) with \(h\in(0,1)\)
admits feasible \(e\), there is a function \(G\) with
\(F(h,q,e)=G(h,q)\).

\emph{Step 2: the negative side depends only on \(q/h\).} Let
\(q\leq h\) and \(r=q/h\in[0,1]\). For feasible \((h,q,e)\),
\(P(EH)=eq\leq eh=P(E)P(H)\), so partial inconsistency gives
\(R(E\Yright H)_P=R(H\Yright E)_P\). The reversed pair has target prior
\(P(E)=e\), source mass \(h\) and posterior \(P(E\mid H)=eq/h=er\), so
\(G(h,rh)=G(e,re)\). Given two priors \(h,h'\in(0,1)\) and the same
\(r\), choose \(e\) with
\(0<e<\min\{(1-h)/(1-rh),(1-h')/(1-rh')\}\); both bounds are positive
and \(e\,rh\leq h\), \(e\,rh'\leq h'\) hold since \(e<1\) and
\(r\leq1\). Then \((h,rh,e)\) and \((h',rh',e)\) are feasible and
\(G(h,rh)=G(e,re)=G(h',rh')\). Hence \(G(h,rh)=\varphi(r)\) for a
function \(\varphi\) on \([0,1]\). If \(r_1<r_2\), take \(h=1/2\),
\(q_i=r_i/2\), and one small common \(e\); the first construction
places both posteriors in one model, and the strict case of POP-D gives
\(\varphi(r_1)<\varphi(r_2)\). So \(\varphi\) is strictly increasing.

\emph{Step 3: the positive side depends only on Z.} Let
\(q_1>h_1\), \(q_2>h_2\), and put \(z_i=(q_i-h_i)/(1-h_i)\in(0,1]\).
Choose \(e\) feasible for both pairs and realize them by the second
construction. The negated targets have priors \(1-h_i\) and posteriors
\(1-q_i<1-h_i\), so by Steps~1 and~2
\(R(E\Yright\neg H_i)_P=\varphi\bigl((1-q_i)/(1-h_i)\bigr)
=\varphi(1-z_i)\). Confirmation complementarity and the strict
monotonicity of \(\varphi\) give
\[
 \begin{aligned}
 G(h_1,q_1)\bowtie G(h_2,q_2)
 &\iff\varphi(1-z_2)\bowtie\varphi(1-z_1)\\
 &\iff 1-z_2\bowtie1-z_1
 \iff z_1\bowtie z_2 .
 \end{aligned}
\]
The equality case shows that \(G(h,q)=\psi(z)\) on the positive side
for a function \(\psi\) on \((0,1]\), every value of which is attained
(take \(h=1/2\), \(q=(1+z)/2\)); the strict cases show that \(\psi\) is
strictly increasing.

\emph{Step 4: one increasing function.} Put \(f(z)=\varphi(1+z)\) for
\(-1\leq z\leq0\) and \(f(z)=\psi(z)\) for \(0<z\leq1\). On the negative
side \(R_Z=(q-h)/h=r-1\), so \(G(h,q)=\varphi(r)=f(R_Z)\); on the
positive side \(G(h,q)=\psi(R_Z)=f(R_Z)\). Each piece is strictly
increasing. For \(q>h\) the first construction places posteriors \(q\)
and \(h\) for one target in one model (the second source independent of
\(H\)), so POP-D gives \(G(h,q)>G(h,h)=\varphi(1)=f(0)\). Hence
\(f(z)>f(0)>f(z')\) whenever \(z>0>z'\), and \(f\) is strictly
increasing on \([-1,1]\), with \(R=f\circ R_Z\) on the comparison domain.

\emph{Item 2.} Let \(0<h<1\). The negated target has prior \(1-h\) and
posterior \(1-q\). If \(q>h\), then \(1-q<1-h\) and
\(R_Z(E\Yright\neg H)_P=((1-q)-(1-h))/(1-h)=-(q-h)/(1-h)\). If
\(q<h\), then \(1-q>1-h\) and
\(R_Z(E\Yright\neg H)_P=((1-q)-(1-h))/(1-(1-h))=-(q-h)/h\). If \(q=h\),
both values are zero. In each case the value is \(-R_Z(E\Yright H)_P\).

\emph{Converse of item 1.} For fixed \(h\), \(R_Z\) is strictly
increasing in \(q\) on each branch and both branches vanish at
\(q=h\), so \(f\circ R_Z\) satisfies POP-D. If
\(P(EH)\leq P(E)P(H)\), then \(q\leq h\) and
\(R_Z(E\Yright H)_P=P(EH)/(P(E)P(H))-1\); also
\(P(E\mid H)=P(EH)/P(H)\leq P(E)\), so
\(R_Z(H\Yright E)_P=P(E\mid H)/P(E)-1\) is the same number, and
partial inconsistency holds for \(f\circ R_Z\). By item~2,
\(R_Z(E\Yright H_1)_P\bowtie R_Z(E\Yright H_2)_P\) exactly when
\(R_Z(E\Yright\neg H_2)_P\bowtie R_Z(E\Yright\neg H_1)_P\), and the
strictly increasing \(f\) preserves both comparisons, so confirmation
complementarity holds.

\emph{Item 3.} With \(P(E\mid H)=eq/h\) and
\(P(E\mid\neg H)=e(1-q)/(1-h)\),
\[
 \begin{aligned}
 P(E\mid H)-P(E)&=\frac{e(q-h)}{h},\\
 P(E\mid H)-P(E\mid\neg H)&=\frac{e(q-h)}{h(1-h)},\\
 R_C(E\Yright H)_P&=eq-eh=e(q-h).
 \end{aligned}
\]
Each is \(e\) times a function of \(h,q\) with the sign of \(q-h\), the
sign of \(R_D\). Take \(h=1/2\), \(q=3/4\), \(e_1=1/4\), \(e_2=1/8\);
both triples are feasible, since \(e_i\leq h/q=2/3\) and
\(e_i\leq(1-h)/(1-q)=2\). The first construction places the two
sources in one model with equal posteriors \(3/4\), so POP-D requires
equal values; each of the three measures gives values in the ratio
\(e_1:e_2=2:1\), with nonzero common factor \(q-h=1/4\).
\end{proof}
\end{annexitem}

\begin{annexitem}{annex:p15-w02b-p15-024-near-certain}{thm:p15-w02b-p15-024-near-certain}{Near-certain entailed evidence}
\begin{proof}
Since \(H\) entails \(D\), \(P(HD)=P(H)=h\). Hence
\(\Delta_P(D,H)=h-dh=h(1-d)>0\), so \(D\) is positive to \(H\), and
\(q=P(H\mid D)=h/d\), which exceeds \(h\) because \(d<1\).

\emph{D and Q.} \(R_D(D\Yright H)_P=h/d-h=h(1-d)/d\), and
\(R_Q(D\Yright H)_P=q/h=1/d\), so \(R_Q-1=(1-d)/d\).

\emph{L.} Since \(0<h<1\), the likelihoods are
\(a=P(D\mid H)=1\) and
\(b=P(D\mid\neg H)=P(D\neg H)/(1-h)=(d-h)/(1-h)\), and \(b>0\) because
\(d>h\). So \(R_L(D\Yright H)_P=(1-h)/(d-h)\) and
\[
 R_L(D\Yright H)_P-1=\frac{(1-h)-(d-h)}{d-h}=\frac{1-d}{d-h}.
\]

\emph{Z.} Because \(q>h\), the first branch of the definition applies:
\[
 R_Z(D\Yright H)_P=\frac{h/d-h}{1-h}=\frac{h(1-d)}{d(1-h)}.
\]

\emph{Limits.} Fix \(h\) and let \(d\to1\) through \((h,1)\). For
\(d\geq(1+h)/2\) we have \(d\geq\tfrac12\) and
\(d-h\geq(1-h)/2\), so the four factors multiplying \(1-d\), namely
\(h/d\), \(1/d\), \(1/(d-h)\) and \(h/(d(1-h))\), are bounded by
\(2h\), \(2\), \(2/(1-h)\) and \(2h/(1-h)\). Hence \(R_D\) and \(R_Z\)
tend to \(0\) and \(R_Q\) and \(R_L\) tend to \(1\), their neutral
values. Since \(\ln\) is continuous at \(1\), \(R_{QL}=\ln R_Q\) and
\(R_{LL}=\ln R_L\) tend to \(0\). For KO, \(0<d,h<1\), so
\(R_{KO}=(a-b)/(a+b)=(1-b)/(1+b)\), and \(b=(d-h)/(1-h)\to1\) gives
\(R_{KO}\to0\).

\emph{Reversed direction.} With source \(H\) and target \(D\),
\(P(H)>0\), \(0<d<1\), and \(P(D\mid H)=1>d\), so the first branch of
Z gives \(R_Z(H\Yright D)_P=(1-d)/(1-d)=1\). Every value of Z is at
most \(1\), since \(P(D\mid H)\leq1\) bounds the first branch by \(1\)
and the second branch is at most \(0\); so \(1\) is the largest value.

\emph{Negated evidence.} \(P(\neg D)=1-d>0\) and \(0<h<1\). The
likelihoods are \(P(\neg D\mid H)=0\), because \(H\) entails \(D\),
and \(P(\neg D\mid\neg H)=(1-d)/(1-h)>0\). So
\(R_L(\neg D\Yright H)_P=0\), and with the convention \(\ln0=-\infty\)
of Definition~\ref{def:p15-r43-measures},
\(R_{LL}(\neg D\Yright H)_P=-\infty\).
\end{proof}
\end{annexitem}

\begin{annexitem}{annex:p15-w02b-p15-024-search}{thm:p15-w02b-p15-024-search}{An unsuccessful search for counterexamples}
\begin{proof}
Write \(t=P(T)\) and \(f=P(F)\). Since \(P(TF)=0\) and \(f>0\),
\(P(T\mid F)=0\), so \(P(\neg T\mid F)=1\); together with \(f>0\) and
\(P(\neg T)=1-t<1\), this is extreme positivity of \(F\) to \(\neg T\),
which is \(F\,\mathrm{exneg}_P\,T\) by
Definition~\ref{def:p15-r43-extreme}.

Also \(P(T\neg F)=P(T)-P(TF)=t\), so
\(\Delta_P(\neg F,T)=t-(1-f)t=tf>0\), and \(\neg F\) is positive to
\(T\). Here \(P(\neg F)=1-f>0\), and
\(q=P(T\mid\neg F)=t/(1-f)\). Hence
\[
 R_D(\neg F\Yright T)_P=\frac{t}{1-f}-t=\frac{tf}{1-f},\qquad
 R_Q(\neg F\Yright T)_P=\frac{q}{t}=\frac1{1-f}.
\]
For L, \(0<t<1\). The likelihood under \(T\) is
\(P(\neg F\mid T)=P(T\neg F)/t=1\). Under \(\neg T\) it is
\(b=P(\neg F\mid\neg T)=1-P(F\mid\neg T)\). If \(P(F\mid\neg T)<1\),
then \(b>0\) and \(R_L(\neg F\Yright T)_P=1/b\). If
\(P(F\mid\neg T)=1\), then \(b=0\) while the likelihood under \(T\) is
\(1\), and Definition~\ref{def:p15-r43-measures} gives
\(R_L(\neg F\Yright T)_P=+\infty\).
\end{proof}
\end{annexitem}

\begin{annexitem}{annex:p15-w02b-p15-024-ko-transform}{thm:p15-w02b-p15-024-ko-transform}{Increasing transforms of L under conditional independence}
\begin{proof}
Write \(h=P(H)\). Since \(0<a_i,b_i<1\) and \(0<h<1\),
\(P(E_i)=ha_i+(1-h)b_i\) lies strictly between zero and one, so
\(R_L\), \(R_{LL}\) and \(R_{KO}\) are given by their main clauses in
Definition~\ref{def:p15-r43-measures}: \(l_i=a_i/b_i\in(0,\infty)\) and
\(k_i=(a_i-b_i)/(a_i+b_i)\).

(1) Dividing numerator and denominator of \(k_i\) by \(b_i>0\) gives
\(k_i=(l_i-1)/(l_i+1)\). For real \(u\),
\(\tanh(u/2)=(e^{u/2}-e^{-u/2})/(e^{u/2}+e^{-u/2})=(e^u-1)/(e^u+1)\);
with \(u=\ln l_i=R_{LL}(E_i\Yright H)_P\) this is \((l_i-1)/(l_i+1)\).
Conditional independence given \(H\) and given \(\neg H\) gives
\(P(E_1E_2\mid H)=a_1a_2\) and \(P(E_1E_2\mid\neg H)=b_1b_2\), both in
\((0,1)\), so \(P(E_1E_2)\in(0,1)\) and
\(R_L(E_1E_2\Yright H)_P=a_1a_2/(b_1b_2)=l_1l_2\). By the first step
applied to \(E_1E_2\),
\(R_{KO}(E_1E_2\Yright H)_P=(l_1l_2-1)/(l_1l_2+1)\). On the other side,
\[
 \frac{k_1+k_2}{1+k_1k_2}
 =\frac{(l_1-1)(l_2+1)+(l_2-1)(l_1+1)}{(l_1+1)(l_2+1)+(l_1-1)(l_2-1)}
 =\frac{2l_1l_2-2}{2l_1l_2+2},
\]
after multiplying numerator and denominator by \((l_1+1)(l_2+1)>0\)
and expanding. The two expressions agree.

(2) The four cells within \(H\) have masses
\(\tfrac12\cdot\tfrac34\cdot\tfrac34\), \(\tfrac12\cdot\tfrac34\cdot\tfrac14\),
\(\tfrac12\cdot\tfrac14\cdot\tfrac34\), \(\tfrac12\cdot\tfrac14\cdot\tfrac14\)
for \(E_1E_2\), \(E_1\neg E_2\), \(\neg E_1E_2\), \(\neg E_1\neg E_2\),
and within \(\neg H\) the same with \(3/4\) and \(1/4\) exchanged;
they are nonnegative and sum to one, and by construction \(E_1\) and
\(E_2\) are independent given \(H\) and given \(\neg H\). Then
\(l_i=(3/4)/(1/4)=3\) and \(k_i=(1/2)/1=1/2\). By (1),
\(R_L(E_1E_2\Yright H)_P=9\) and
\(R_{KO}(E_1E_2\Yright H)_P=(9-1)/(9+1)=4/5\); directly,
\((9/16-1/16)/(9/16+1/16)=4/5\). Finally \(k_1+k_2=1\neq4/5\) and
\(k_1k_2=1/4\neq4/5\).

(3) First, every pair \((x,y)\in(0,\infty)^2\) occurs as \((l_1,l_2)\)
in a model of the stated kind. Take \(h=1/2\); for \(l\in\{x,y\}\) put
\(b(l)=1/(2\max\{1,l\})\) and \(a(l)=l\,b(l)\), so that
\(0<a(l),b(l)\leq1/2\); give the eight cells
\(H^{\pm}E_1^{\pm}E_2^{\pm}\) the product masses as in (2). Then
\(l_1=x\), \(l_2=y\), conditional independence holds, and by (1)
\(R_L(E_1E_2\Yright H)_P=xy\).

If \(f=c\ln\) with \(c>0\), then \(f(l_1l_2)=f(l_1)+f(l_2)\) by the
logarithm law, and \(f\) is strictly increasing. If \(f(x)=x^c\) with
\(c>0\), then \(f(l_1l_2)=f(l_1)f(l_2)\) and \(f\) is strictly
increasing. These are the converse directions.

Suppose the sum rule holds in every model. By the realization step,
\(f(xy)=f(x)+f(y)\) for all \(x,y>0\). Put \(g(t)=f(e^t)\) for real
\(t\); then \(g(s+t)=g(s)+g(t)\) and \(g\) is strictly increasing.
From \(g(0)=g(0)+g(0)\), \(g(0)=0\). Induction gives \(g(nt)=ng(t)\)
for positive integers \(n\); with \(g(t)+g(-t)=g(0)=0\) this extends
to all integers. For a rational \(r=m/n\) with \(n>0\),
\(ng(r)=g(m)=mg(1)\), so \(g(r)=rg(1)\). Put \(c=g(1)\); then
\(c>g(0)=0\). For real \(t\) and rationals \(r<t<r'\), monotonicity
gives \(rc<g(t)<r'c\); letting \(r\) and \(r'\) tend to \(t\) yields
\(g(t)=ct\). Hence \(f(x)=g(\ln x)=c\ln x\).

Suppose the product rule holds in every model, so that
\(f(xy)=f(x)f(y)\) for all \(x,y>0\). Then
\(f(x)=f(\sqrt x)^2\geq0\). If \(f(x_0)=0\) for some \(x_0\), then
\(f(x)=f(x_0)f(x/x_0)=0\) for every \(x\), contradicting strict
increase; so \(f>0\). The function \(\ln f\) is strictly increasing and
satisfies \(\ln f(xy)=\ln f(x)+\ln f(y)\), so by the previous paragraph
\(\ln f=c\ln\) with \(c>0\), that is, \(f(x)=x^c\).

KO is \(f(R_L)\) with \(f(x)=(x-1)/(x+1)\), which is strictly
increasing on \((0,\infty)\) and is neither of these forms, as (2)
shows directly.
\end{proof}
\end{annexitem}

\begin{annexitem}{annex:p15-w02b-p15-024-mdhd-ratios}{thm:p15-w02b-p15-024-mdhd-ratios}{D, Q and L on a refuted hypothesis and its disjoint weakening}
\begin{proof}
Write \(e=P(E)\). By Definition~\ref{def:p15-r43-extreme},
\(E\,\mathrm{exneg}_P\,H\) means \(e>0\), \(P(\neg H)<1\) and
\(P(\neg H\mid E)=1\); so \(h>0\) and \(P(EH)=0\). Independence gives
\(P(EG)=eg\), and disjointness of \(H\) and \(G\) gives
\(P(H\vee G)=h+g<1\) and
\(P(E(H\vee G))=P(EH)+P(EG)=eg\). Also \(g\leq1-h<1\).

(1) \(P(H\mid E)=0\), so \(R_D(E\Yright H)_P=-h\). And
\(P(H\vee G\mid E)=eg/e=g\), so
\(R_D(E\Yright H\vee G)_P=g-(h+g)=-h\).

(2) \(R_Q(E\Yright H)_P=0/h=0\) and
\(R_Q(E\Yright H\vee G)_P=g/(h+g)\), positive since \(g>0\).

(3) For \(H\): \(0<h<1\), \(P(E\mid H)=0\) and
\(P(E\mid\neg H)=e/(1-h)>0\), so \(R_L(E\Yright H)_P=0\). For
\(H\vee G\), \(0<h+g<1\);
\(P(E\mid H\vee G)=eg/(h+g)\); and
\(P(E\neg(H\vee G))=e-eg=e(1-g)>0\), so
\(P(E\mid\neg(H\vee G))=e(1-g)/(1-h-g)>0\). The ratio is
\(g(1-h-g)/((h+g)(1-g))\), positive because each factor is.

(4) \(P(H\vee G\mid E)=g<h+g=P(H\vee G)\), so \(E\) is negative to
\(H\vee G\). It is not extremely negative, since that would require
\(P(\neg(H\vee G)\mid E)=1\), whereas this probability is \(1-g<1\).

For the consequence: the D values are equal, while the Q and the L
values satisfy a strict inequality in favour of \(H\vee G\), which every
strictly increasing transform preserves. Ordinal equivalence of D with
one of these measures would require the comparison
\(R(E\Yright H\vee G)_P\leq R(E\Yright H)_P\) to hold for both
measures or fail for both; it holds for D and fails for the other. In
the card model \(E\) and \(G\) are independent, since
\(P(EG)=6/52=(12/52)(1/2)\); \(P(HG)=0\); and \(E\) excludes \(H\).
Substituting \(h=1/52\) and \(g=1/2\) gives
\(g/(h+g)=(1/2)/(27/52)=26/27\) and
\(g(1-h-g)/((h+g)(1-g))=(1/2)(25/52)/((27/52)(1/2))=25/27\).
\end{proof}
\end{annexitem}

\begin{annexitem}{annex:p15-w02b-p15-024-diagnostic-additivity}{thm:p15-w02b-p15-024-diagnostic-additivity}{Additivity of screened symptoms}
\begin{proof}
Write \(h=P(B)\in(0,1)\). For any event \(X\) with
\(x_1=P(X\mid B)\) and \(x_0=P(X\mid\neg B)\) not both zero,
\(P(X)=hx_1+(1-h)x_0>0\) and
\[
 P(B\mid X)-h=\frac{hx_1}{hx_1+(1-h)x_0}-h
 =\frac{h(1-h)(x_1-x_0)}{P(X)},
\]
so \(X\) is positive to \(B\) exactly when \(x_1>x_0\).

(1) Conditional independence given \(B\) and given \(\neg B\) gives
\(P(AC\mid B)=a_1a_2\) and \(P(AC\mid\neg B)=b_1b_2\), both positive.
Hence \(R_L(AC\Yright B)_P=a_1a_2/(b_1b_2)=(a_1/b_1)(a_2/b_2)\), the
product of the two likelihood ratios, all finite and positive; the
logarithm of the product is the sum of the logarithms.

(2) Positivity of \(A\) and of \(C\) to \(B\) gives, by the first
paragraph, \(a_1>b_1\geq0\) and \(a_2>b_2\geq0\). Then
\(a_1a_2>b_1a_2\geq b_1b_2\), since \(a_2>0\) and \(b_1\geq0\); so the
likelihoods of \(AC\) satisfy the same strict inequality, and \(AC\) is
positive to \(B\). If \(b_1,b_2>0\), then \(l_i=a_i/b_i>1\), so
\(\ln l_i>0\) and \(\ln(l_1l_2)=\ln l_1+\ln l_2\) exceeds each
summand, which is the stated chain.

(3) By (2), under \((A\perp C\mid\pm B)_P\) and positivity of \(A\) and
\(C\) to \(B\), \(AC\) is positive to \(B\). So if \(AC\) is irrelevant
or negative to \(B\) while \(A\) and \(C\) are positive to it, the
premise \((A\perp C\mid\pm B)_P\) is false.
\end{proof}
\end{annexitem}

\begin{annexitem}{annex:p15-w02b-p15-024-exneg-disjuncts}{thm:p15-w02b-p15-024-exneg-disjuncts}{Disjoint disjuncts}
\begin{proof}
By Definition~\ref{def:p15-r43-extreme}, \(A\,\mathrm{exneg}_P\,B\)
means \(P(A)>0\), \(P(\neg B)<1\) and \(P(\neg B\mid A)=1\).

(1) From \(P(\neg B)<1\), \(P(B)>0\). From \(P(\neg B\mid A)=1\),
\(P(A\neg B)=P(A)\), so \(P(AB)=P(A)-P(A\neg B)=0\). For the converse
relation: \(P(B)>0\); \(P(\neg A)=1-P(A)<1\); and
\(P(\neg A\mid B)=(P(B)-P(AB))/P(B)=1\). So \(B\,\mathrm{exneg}_P\,A\).

(2) If \(A\,\mathrm{exneg}_P\,B\), then \(P(AB)=0\) by (1).
Conversely, if \(P(A),P(B)>0\) and \(P(AB)=0\), then
\(P(\neg B)=1-P(B)<1\) and
\(P(\neg B\mid A)=(P(A)-P(AB))/P(A)=1\), so \(A\,\mathrm{exneg}_P\,B\);
by (1) also \(B\,\mathrm{exneg}_P\,A\).

(3) Since \(P(A)=0\), \(P(X(A\vee B))=P(XB)\) for every event \(X\),
and in particular \(P(A\vee B)=P(B)\). Each of D, Q and L, including
its completion clauses, is computed from the probabilities of the
source, of the target and of their Boolean combinations, so
\(R_T(A\vee B\Yright C)_P=R_T(B\Yright C)_P\) for \(T=D,Q,L\). The
D and Q weight is \(P(A)/(P(A)+P(B))=0\). For L the weight is
\(P(A\mid\neg C)/(P(A\mid\neg C)+P(B\mid\neg C))\); when
\(0<P(C)<1\), \(P(A\mid\neg C)=0\), and when \(P(B\mid\neg C)>0\) the
denominator is positive and the weight is \(0\).
\end{proof}
\end{annexitem}
